\section{Related Work}\label{sec:related}

\noindent\textbf{MPPI and Deployable Local Controllers.} Nav2 MPPI \citep{nav2,mppi_orig}, TEB \citep{teb}, DWB, a Nav2 planner of the dynamic-window family \citep{dwa}, and Regulated Pure Pursuit \citep{rpp} are the standard deployable local controllers. MPPI's critic interface (\code{mppi::critics::CriticFunction}) exposes the per-rollout trajectories, the reference path, the goal and the furthest-reached path index, which are exactly the signals needed for entrapment detection, so a custom critic can be inserted without forking the optimizer. \sname{} uses this interface for escape and adds an output-side safety filter, instead of replacing the optimizer. Around the controller, Nav2 already handles stalls at the stack level: the controller server runs a progress checker, by default \code{SimpleProgressChecker}, which reports a failure when the robot has not moved a set radius within a set time allowance, and the default behaviour tree answers with costmap clearing, spinning, waiting and backing up \citep{nav2_progress,nav2_bt}. \sname's detector differs in that it runs inside the controller on progress along the path and steers the escape instead of triggering generic recoveries; the two mechanisms can coexist.

\noindent\textbf{Local-Minima Escape.} Most escape methods for MPPI are always-on changes of the sampling distribution: log-MPPI \citep{log_mppi}, Tsallis-divergence MPC \citep{tsallis_mppi}, Stein-variational guidance \citep{svg_mppi}, ancillary-controller biasing \citep{biased_mppi}, and learned sampling distributions \citep{flow_mppi}. Others add guidance. GP-MPPI learns a sparse-Gaussian-process model of the navigable space around the robot and recommends a free-space subgoal to the local MPPI planner, which lets the robot escape local minima without a global map \citep{gp_mppi}. Fuke et al.\ add a repulsive artificial potential to MPPI to avoid local minima without global path guidance \citep{drpa_sii}, and their DRPA-MPPI detects potential entrapment on the predicted trajectories and only then switches to a cost with repulsion away from the local minimum \citep{drpa_mppi}. DRPA-MPPI is the closest detect-and-switch method and GP-MPPI the closest subgoal method; neither has a CBF or other formal safety layer, and neither is Nav2-integrated. \sname{} combines both ideas in a simpler form, a path-progress stall detector gating a classical repulsive potential \citep{khatib_apf} and a ray-cast gap bearing, and its safety filter runs alongside the escape rather than gating it. Gap-based local planners are closer still to the overall architecture. Potential Gap combines gap-based local navigation with artificial potential fields and restores provable collision-free properties for realistic robot models \citep{potential_gap}; Safer Gap adds nonlinear model predictive control with keyhole zeroing-barrier-function constraints and a point-wise optimization-based safety filter as the final layer \citep{safer_gap}; and Dynamic Gap extends provably collision-free gap-based navigation to dynamic environments by modelling how gaps evolve over time \citep{dynamic_gap}. Relative to this line, \sname{} uses a gap only as a gated escape heuristic inside MPPI, and its barrier covers only tracked dynamic obstacles.

\noindent\textbf{Escape and Safety Together.} Several works treat entrapment and safety jointly. Soloperto et al.\ optimize an exploration trajectory together with a safe backup trajectory in a model predictive controller and prove convergence, feasibility and constraint satisfaction under partially unknown constraints \citep{soloperto}. The modulated CBF-QP of Xue and Figueroa reduces or eliminates, in dynamic environments, the spurious equilibria that CBF-QP filters create \citep{reis_cbf_equilibria}, and was tested on Ridgeback and Fetch robots \citep{mcbf_qp}. Concurrent work adds learned predictive barriers for moving concave obstacles to this filter \citep{xue_proactive}, and pairs a mixed-integer MPC planner with a Minkowski-difference CBF filter to mitigate the local-minimum behaviour of reactive CBF navigation in U-shaped and maze-like environments \citep{chen_milp_cbf}. These methods mitigate or avoid minima inside the planner or the barrier itself, whereas \sname{} leaves the barrier unchanged and escapes through the MPPI cost. None of them is packaged as a Nav2 controller.

\noindent\textbf{CBF--MPPI Safety.} Three fusion styles exist as research code: CBF conditions in the rollout cost or the sampling (Shield-MPPI \citep{shield_mppi}, which adds a local repair step, and BR-MPPI \citep{br_mppi}, which imposes the CBF condition as an equality with a parametric linear class-$\mathcal{K}$ function, augments the state with its parameter and lets MPPI optimize the parameter's rate as an extra input), an output CBF-QP filter on the MPPI command (SCBF-MPPI \citep{scbf_mppi}, built on the CBFKit toolbox \citep{cbfkit}), and per-rollout safety filtering that makes every sample provably safe (GS-MPPI \citep{gs_mppi}, DualGuard \citep{dualguard}, which also filters the final command). For dynamic obstacles, DPCBF \citep{dpcbf} uses a parabolic safe set that is less conservative than collision-cone barriers but is QP-only and not embedded in MPPI. Our filter builds on the classical CBF-QP \citep{ames_cbf_tac,ames_cbf_ecc}. None of these works ships as a Nav2 controller (CBFKit builds a standalone node), and none detects entrapment or adds an escape mechanism, although GS-MPPI counters the myopic behaviour often associated with CBFs through long-term performance terms \citep{gs_mppi}. Among Nav2 integrations of a CBF that we could verify, the official documentation describes inserting the third-party 3Laws Supervisor, which its vendor describes as based on CBFs, between the velocity smoother and the collision monitor \citep{nav2_3laws,threelaws_pr}; it acts on the controller's output as a separate node, and neither source describes a local-minimum escape mechanism. \sname{} instead places its filter inside the controller and couples it to the escape layer through its gain.

\noindent\textbf{Online Adaptation of the CBF Gain.} Time-varying and adapted class-$\mathcal{K}$ terms are well studied. Adaptive CBFs multiply the class-$\mathcal{K}$ functions by time-varying penalty functions and prove forward invariance of the safe set \citep{ada_cbf}; optimal-decay CBF-QPs optimize the decay rate point-wise in time to keep the QP feasible \citep{optdecay_cbf}; and rate-tunable CBFs adapt the class-$\mathcal{K}$ functions online and derive conditions under which multiple constraints can share the control input \citep{rt_cbf}. Closest to our trigger, Kim, Kee and Panagou adapt the parameters of input-constrained CBFs online with a learned ensemble model and select, among the candidates that pass their uncertainty and risk checks, the class-$\mathcal{K}$ parameters with the minimum predicted deadlock time \citep{kim_adaptive_cbf}; BR-MPPI lets MPPI optimize the class-$\mathcal{K}$ parameter directly \citep{br_mppi}. Concurrent work switches nominal objectives when CBF and control-Lyapunov constraints conflict, to reduce slowdowns and deadlocks \citep{walia_conflict}, and schedules the gain of a frontier barrier on an occupancy grid by how well the map is explored \citep{paudel_dualbarrier}. \sname's coordinator is a rule-based, two-level instance of this idea, triggered by the entrapment flag and overridden by time-to-collision, and its invariance property (\cref{prop:invariance}) follows the same argument as the relative-degree-one case of the adaptive-CBF result \citep[Thm.~3]{ada_cbf}, there stated for Lipschitz controllers.

\noindent\textbf{Conformal Prediction for Safe Control.} The released controller also carries an online conformal margin that inflates the barrier radius under prediction uncertainty, drawing on conformal prediction for safe planning and control \citep{acp_yang,conformal_hri,lindemann,dixit} and on uncertainty-aware predictive barriers \citep{ua_pcbf}. This margin is follow-up work and is \emph{not} a contribution of this paper; every result reported here is isolated from it (\cref{sec:discussion}).

\begin{table}[t]
\centering
\caption{\textbf{Positioning.} (a)~ships as a Nav2 controller plugin; (b)~online local-minimum escape or avoidance; (c)~an explicit safety constraint, filter or backup enforced at run time (for DualGuard, Hamilton--Jacobi reachability; for Soloperto et al., a safe backup trajectory; for Safer Gap, zeroing barrier functions); (d)~online adaptation of the CBF class-$\mathcal{K}$ gain. \cmark{} = holds; \pmark{} = holds in part: stock Nav2 detects stalls at the stack level and runs generic recoveries rather than a steered escape; the 3Laws Supervisor is integrated with Nav2 as a separate node after the controller, not as a controller plugin; the gap-based planners guide the robot through free-space gaps rather than escaping a detected minimum, and the provable collision-free properties of Potential Gap and Dynamic Gap follow from the planners' gap construction rather than from an explicit run-time constraint or filter; Kim et al.\ reduce predicted deadlock through class-$\mathcal{K}$ parameter selection rather than escaping a detected minimum. For the 3Laws Supervisor, (b) is marked absent because neither cited source describes a local-minimum escape mechanism. $^{\dagger}$Concurrent work (posted in 2026). For \sname, (c) is marked in part because its barrier covers only tracked dynamic obstacles, at a look-ahead point, and its invariance argument is conditional (\cref{prop:invariance}); static structure is handled only by MPPI costs. Its (b) is demonstrated only in the Python 2D mirror, not in a live Nav2 run. No reviewed system holds all four; several hold two. The review behind this table (June 2026, extended in September 2026) covered the Nav2 MPPI critic ecosystem, the Nav2 documentation, the public issue trackers of these projects, and the works cited here; it is not exhaustive.}
\label{tab:positioning}
\small
\setlength{\tabcolsep}{4pt}
\begin{tabularx}{\linewidth}{>{\raggedright\arraybackslash}p{0.34\linewidth}>{\raggedright\arraybackslash}Xcccc}
\toprule
System & Reference & (a) & (b) & (c) & (d) \\
\midrule
Nav2 MPPI with progress checker and recoveries & \citet{nav2,nav2_progress,nav2_bt} & \cmark & \pmark & \xmark & \xmark \\
3Laws Supervisor in Nav2 & \citet{nav2_3laws,threelaws_pr} & \pmark & \xmark & \cmark & \xmark \\
DRPA-MPPI, GP-MPPI & \citet{drpa_mppi,gp_mppi} & \xmark & \cmark & \xmark & \xmark \\
Dynamic Gap, Potential Gap & \citet{dynamic_gap,potential_gap} & \xmark & \pmark & \pmark & \xmark \\
Safer Gap & \citet{safer_gap} & \xmark & \pmark & \cmark & \xmark \\
Backup-trajectory MPC & \citet{soloperto} & \xmark & \cmark & \cmark & \xmark \\
MILP-MPC with CBF$^{\dagger}$; MCBF-QP; predictive MCBF$^{\dagger}$ & \citet{chen_milp_cbf,mcbf_qp,xue_proactive} & \xmark & \cmark & \cmark & \xmark \\
DualGuard, SCBF-MPPI, GS-MPPI, Shield-MPPI & \citet{dualguard,scbf_mppi,gs_mppi,shield_mppi} & \xmark & \xmark & \cmark & \xmark \\
CBFKit, DPCBF & \citet{cbfkit,dpcbf} & \xmark & \xmark & \cmark & \xmark \\
BR-MPPI & \citet{br_mppi} & \xmark & \xmark & \cmark & \cmark \\
Rate-tunable, adaptive, optimal-decay CBFs & \citet{rt_cbf,ada_cbf,optdecay_cbf} & \xmark & \xmark & \cmark & \cmark \\
Deadlock-aware parameter adaptation & \citet{kim_adaptive_cbf} & \xmark & \pmark & \cmark & \cmark \\
\rowcolor{black!6}\textbf{\sname{} (ours)} & this paper & \cmark & \cmark & \pmark & \cmark \\
\bottomrule
\end{tabularx}
\end{table}


\noindent\textbf{Positioning.} \Cref{tab:positioning} places the reviewed systems against four factors: shipping as a Nav2 controller plugin, online local-minimum escape or avoidance, an explicit run-time safety constraint, filter or backup, and online adaptation of the class-$\mathcal{K}$ gain. No reviewed system holds all four, but several hold two, fully or in part: stock Nav2 pairs its controller plugins with stack-level stall recovery, gap-based planners, backup-trajectory MPC and modulated CBFs combine local-minimum handling with a safety layer, and adaptive-gain CBFs and BR-MPPI pair a barrier with gain adaptation, with deadlock-aware adaptation also holding the second factor in part. The delta of \sname{} is therefore the conjunction inside one Nav2 controller plugin, not any single element. Its own marks are not all full: its barrier covers only tracked dynamic obstacles, and its escape is demonstrated only in the Python 2D mirror. Its gain rule differs from prior gain adaptation mainly in its trigger, and, as \cref{sec:ablation} shows, we could not measure an outcome benefit from it.
